\documentclass[12pt]{article}
\usepackage{amsmath,amssymb,amsfonts}
\usepackage[margin=1in]{geometry}

\begin{document}

\begin{center}
{\Large\bfseries Chapter 10, Problem 13}
\end{center}

\noindent\textbf{Problem.} Show that if $C_i$ is a class containing $S_1, S_2, \ldots, S_g$ and $C_i^{-1}$ is the collection of elements consisting of $S_1^{-1}, S_2^{-1}, \ldots, S_g^{-1}$, then $C_i^{-1}$ is also a class. Show that for a general irreducible representation, i.e., one which is not necessarily unitary,
\[
\sum_g g|\chi_i^{(\nu)}|^2 = g\delta_{\nu\nu}, \qquad a_\nu = \frac{1}{g}\sum_g g_i\chi_i^{(\nu)*}\chi_i,
\]
where $a_\nu$ is the number of times the $\nu$th irreducible representation occurs in the representation (not necessarily irreducible) whose characters are $\chi_i$.

\bigskip
\noindent\textbf{Solution.}

\noindent\textbf{Part 1: $C_i^{-1}$ is a class.}

Let $C_i = \{S_1, S_2, \ldots, S_g\}$ be a conjugacy class. We need to show $C_i^{-1} = \{S_1^{-1}, S_2^{-1}, \ldots, S_g^{-1}\}$ is also a conjugacy class.

First, $C_i^{-1}$ is contained in a single conjugacy class: if $S_j \in C_i$, then for any $T \in G$:
\[
T S_j^{-1} T^{-1} = (T S_j^{-1} T^{-1}) = (T^{-1})^{-1} S_j^{-1} (T^{-1}) .
\]
But also, since $S_j$ and $S_k$ are conjugate (both in $C_i$), there exists $U$ with $S_k = U S_j U^{-1}$. Then:
\[
S_k^{-1} = (U S_j U^{-1})^{-1} = U S_j^{-1} U^{-1}.
\]
So $S_k^{-1}$ is conjugate to $S_j^{-1}$. This means all elements of $C_i^{-1}$ are mutually conjugate.

Now suppose $T$ is conjugate to $S_1^{-1}$, i.e., $T = U S_1^{-1} U^{-1}$ for some $U$. Then $T^{-1} = U S_1 U^{-1} \in C_i$, so $T^{-1} = S_k$ for some $k$, hence $T = S_k^{-1} \in C_i^{-1}$. Therefore $C_i^{-1}$ contains exactly all elements conjugate to any of its members, so it is a complete conjugacy class. $\square$

\medskip
\noindent\textbf{Part 2: Orthogonality for general (not necessarily unitary) irreducible representations.}

For a finite group, every representation is equivalent to a unitary representation via a similarity transformation $S$: if $D$ is any representation, there exists an invertible $S$ such that $D'(g) = S D(g) S^{-1}$ is unitary for all $g$.

The characters are invariant under similarity transformations:
\[
\chi^{(\nu)}(g) = \operatorname{Tr} D^{(\nu)}(g) = \operatorname{Tr} D'^{(\nu)}(g).
\]

For unitary representations, we have the standard orthogonality relation (Eq.~10.49 in the text):
\[
\sum_{i} g_i \chi_i^{(\mu)*} \chi_i^{(\nu)} = g\,\delta_{\mu\nu},
\]
where the sum is over classes, $g_i$ is the number of elements in class $i$, and $g = |G|$.

Since the characters of the general irreducible representation are the same as those of the equivalent unitary representation, the orthogonality relation holds unchanged:
\[
\boxed{\sum_{i} g_i \chi_i^{(\mu)*} \chi_i^{(\nu)} = g\,\delta_{\mu\nu}.}
\]

Similarly, for the decomposition of an arbitrary representation with character $\chi$ into irreducible representations, the number of times the $\nu$th irreducible representation appears is:
\[
\boxed{a_\nu = \frac{1}{g}\sum_{i} g_i \chi_i^{(\nu)*} \chi_i.}
\]

This follows because characters depend only on the equivalence class of a representation, and every representation of a finite group is equivalent to a unitary one. $\square$

\end{document}
